\section{Polynomial Agreement Certificates}\label{sec:polynomial-agreement}

Fix the RS pair code \(\Cpair\) with distinct evaluation points \(a_i\),
block length \(n\), and degree bound \(1\le k\le n\). Write the received pair
as \(w_i=(w_{i0},w_{i1})\). The quantity certified here is a lower bound on
the whole-pair error count
\(e(w)=\min_{c\in\Cpair}d_H(w,c)\), before the scalar challenge.

\subsection{Agreement and the normalized coefficient system}

\begin{definition}[Polynomial agreement certificate]
A certificate consists of polynomials \(h_0,h_1,R\in\F[X]\), a multiplier
budget \(t\ge0\), and a response degree \(D=\deg R\) satisfying
\begin{equation}\label{eq:poly-agreement-eval}
  \deg h_j\le t,\qquad k+t\le D<n,\qquad
  R(a_i)=h_0(a_i)w_{i0}+h_1(a_i)w_{i1}\quad(i\in I).
\end{equation}
Its certified count is \(n-D\). Write \(b(w)\) for the greatest certified
count, with value zero when the certificate set is empty.
\end{definition}

\begin{proposition}[Soundness]\label{prop:poly-agreement-sound}
Every certificate bounds whole-pair agreement with any valid pair by \(D\)
sites. Consequently \(b(w)\le e(w)\).
\end{proposition}
\begin{proof}
Let \(f_0,f_1\) of degrees below \(k\) encode the valid pair. At every
matching site the polynomial \(R-h_0f_0-h_1f_1\) vanishes by
\cref{eq:poly-agreement-eval}. Its degree is exactly \(D\) because the
subtracted polynomial has degree below \(k+t\le D\). The distinct
evaluation points give at most \(D\) roots.
\end{proof}

\begin{proposition}[Exact feasibility and the degree boundary]
\label{prop:poly-agreement-boundary}
Fix \(k+t\le D<n\). A certificate of budget \(t\) and degree \(D\) exists
if and only if the following system, in the coefficients \(h_{js},r_s\),
is solvable over \(\F\).
\begin{equation}\label{eq:poly-agreement-matrix}
  \sum_{j=0}^1 w_{ij}\sum_{s=0}^{t}a_i^s h_{js}
       -\sum_{s=0}^{D-1} a_i^s r_s=a_i^D
       \qquad(i\in I).
\end{equation}
A nonempty solution set for fixed \(w,t,D\) is a translate of the kernel
of its coefficient matrix. With maximal budget \(t=D-k\) at each degree,
the feasible degrees form an upper interval in \(\{k,\ldots,n-1\}\).
\end{proposition}
\begin{proof}
Divide a certificate's three polynomials by the leading coefficient of
\(R\). Its normalized response is \(X^D+\sum_{s<D}r_sX^s\), and evaluation
gives \cref{eq:poly-agreement-matrix}. Conversely, those normalized
coefficients and \(h_j=\sum_{s\le t}h_{js}X^s\) give a certificate.
For a fixed right-hand side, the difference of any two solutions lies in
the kernel, and adding a kernel element to a solution preserves it.

For fixed \(D\), budget \(D-k\) admits every witness at a smaller budget.
Given a witness at \(D\), multiplication of all three polynomials by
\(X^{E-D}\) gives a witness at any \(D\le E<n\) within budget \(E-k\).
The evaluation equations are preserved even at an evaluation point equal
to zero. The response stays monic and has degree \(E\).
\end{proof}

Finite-field elimination and binary search over the degree interval find
the first feasible degree using only \(w\) and the evaluation points. For
the systematic modular profile in Section~7, the same certificate on the
extracted pair bounds the cochain distance by \(d_C\ge(n-D)/(2n)\), using
\(D_{\mathrm{RS}}\le2d_C\).

\subsection{Closure of the received-word locus}

Now vary \(w\) over an algebraic closure \(K\) of \(\F\) and keep the same
evaluation points, \(t\), and \(D\). Let \(J\subseteq K[w,h,r]\) be the
ideal generated by the row equations in \cref{eq:poly-agreement-matrix},
with their right-hand sides moved to the left. Projection of \(V(J)\) to
the \(w\) coordinates gives the feasible locus. Its Zariski closure is
\(V(J\cap K[w])\) by the closure theorem for elimination
\cite[Theorem~6.17.1]{setiabrata2019computationalAlgebra}.

This closure can add a received pair with empty solution fiber. For
\(\lambda\ne0\), the pair \(w_{\lambda,i}=(\lambda a_i^D,0)\) has multipliers
\((1/\lambda,0)\) and response \(X^D\). Its specialization at \(\lambda=0\)
has empty fiber because a degree-\(D\) monic polynomial cannot vanish at all
\(n>D\) sites. Thus the affine kernel description is a statement about
fixed-data fibers, and elimination gives the closure across received data.

At four sites with \(t=0\) and \(k=D=1\), write \(a=(a_i)_i\) and
\(\mathbf 1=(1)_i\). The closure has the explicit equation
\(\det[\mathbf 1,a,w_0,w_1]=0\). A solution gives a dependence among these
four columns. Conversely, when the determinant is zero, some coefficients
\((\alpha,\beta)\ne(0,0)\) satisfy
\(\alpha w_0+\beta w_1\in\operatorname{span}(\mathbf 1,a)\).
If the combination's \(a\) coefficient is zero,
add \(\lambda a\) to a received column whose coefficient in the combination
is nonzero. For \(\lambda\ne0\) the perturbed pair admits a certificate,
and the original pair is in the closure by specialization.

\subsection{Completeness and worst-case sharpness}

\begin{theorem}[Completeness guarantee and sharpness on nonzero domains]
\label{thm:poly-agreement-complete}
Let \(d_{\min}=n-k+1\) and write \(e=e(w)\). Then
\[
  \min(e,d_{\min}-e)\le b(w)\le e.
\]
In particular \(b(w)=e\) when \(2e\le d_{\min}\), including equality.
If all evaluation points are nonzero, then for each error count
\(0\le e\le n-k\) there is a received pair with true count \(e\) and
certified count exactly \(\min(e,d_{\min}-e)\).
\end{theorem}
\begin{proof}
The upper bound follows from \cref{prop:poly-agreement-sound} and implies
\(b(w)=0\) when \(e=0\). For \(e>0\), choose a nearest pair encoded by
\(f_0,f_1\), and let \(B\) be its agreement set of size \(M=n-e\). Set
\[
  \begin{gathered}
    D=\max(M,k+e-1),\\
    V_B(X)=\prod_{i\in B}(X-a_i),\qquad Q(X)=X^{D-M}V_B(X).
  \end{gathered}
\]
Interpolation on any \(k\) sites gives \(e\le n-k\), so \(D<n\).
At each of the \(e\) mismatching sites the error symbol is nonzero, and we
choose two field values whose dot product with that error symbol equals
\(Q(a_i)\). Interpolating the two lists on those sites gives \(h_0,h_1\) of
degrees at most \(e-1\). Set \(R=Q+h_0f_0+h_1f_1\). The degree of the
added polynomial is below \(k+e-1\le D\), so \(R\) is monic of degree
\(D\). The prescribed products on mismatching sites and the vanishing of
\(Q\) on matching sites verify \cref{eq:poly-agreement-eval}. The count is
\(n-D=\min(e,d_{\min}-e)\). This argument proves existence using a nearest
pair. The search in \cref{prop:poly-agreement-boundary} needs only the
received data.

For sharpness choose any \(B\) of size \(M=n-e\), set \(s=M-k+1\), and
take \(w_i=(V_B(a_i)/a_i^s,0)\). The zero pair agrees on \(B\). Any pair
encoded by \(f_0,f_1\) can agree only at roots of \(V_B-X^sf_0\). This
polynomial has degree at most \(M\) and nonzero constant term \(V_B(0)\),
since every \(a_i\ne0\) and \(s\ge1\). Hence the true count is \(e\).
For any certificate of degree \(D\), the polynomials \(X^sR\) and
\(h_0V_B\) agree at all \(n\) sites. The former has degree \(D+s\) and
the latter has degree at most \(D+s-1\), by \(k+t\le D\). Their difference
is nonzero, so \(D+s\ge n\). Every certified count is at most
\(s=d_{\min}-e\). Together with soundness and the existence bound this
gives the claimed equality.
\end{proof}

For example, on \(\mathbb F_5^\times\) with \(k=1\), choose \(B\) to be
the single site at \(1\). The true count is three and the certified count
is one. The same construction on the canonical 4,096-site domain with
\(|B|=1025\) has true count \(3071\) and certified count two.

On that same domain, another pair with true count \(3071\) admits an exact
certificate. Set \(R=V_B\), \(P_0=X^{3069}R\), and \(P_1=XP_0+R\), and
receive their evaluations. The zero pair agrees on \(B\), so \(e\le3071\).
The multipliers \((-X,1)\) of budget one return \(R\) of degree
\(1025=k+1\), hence \(b=e=3071\). The interpolants \(P_0,P_1\) have degrees
\(4094,4095\), so every nonzero constant combination has degree at least
\(4094\). The best count certified with constant multipliers is two.

\subsection{An application to Euclidean tensor prices}

Place the evaluation sites and a hub in Euclidean space, and let \(S\) be
the set of all their positions, assumed to have positive diameter. Write
\(\ell_{(1)}\ge\cdots\ge\ell_{(n)}\) for the ordered hub-to-site lengths.
Suppose \(H\) is a compatible price system with the wire, product,
restriction, and direct-sum laws of \cite{openai2026euclideanTensorPrices},
with \(W(S)\) its mean-width normalization. A wire with index dimension
\(2^m\) costs \(m\) times its length. For each codeword pair \(c\), take the
product of these wires from the hub to the sites where \(c_i=w_i\), with
a basis tag for \(c\) in a separate factor at the hub. At each unused wire
slot, put a fixed basis vector at both endpoints so that the summands share
the same local spaces. Take their ordinary sum \(T_m\). These product
factors leave the summand price unchanged. Selecting any tag recovers its
summand by local restriction, and the ordinary sum is a local restriction
of the direct sum. Therefore
\[
  m\max_{c\in\Cpair}\sum_{i:c_i=w_i}\ell_i
    \ \le\ H(T_m;S)\ \le\
  m\max_{c\in\Cpair}\sum_{i:c_i=w_i}\ell_i
    +W(S)\log_2|\Cpair|.
\]
By \cref{prop:poly-agreement-sound}, a degree-\(D\) certificate replaces
the maximum in the upper bound by \(\sum_{j=1}^D\ell_{(j)}\).
For unit lengths the limit \(H(T_m;S)/m\) is \(n-e(w)\).
For general lengths the corresponding limit is the maximal weighted
agreement. The Lean results cover the finite certificate arguments in
the first and third subsections. Elimination for the received-word locus
and the application of the Euclidean price theorem are mathematical
arguments in this manuscript.
